\documentclass[11pt,a4paper]{article}

\usepackage[T1]{fontenc}
\usepackage{lmodern}
\usepackage[utf8]{inputenc}
\usepackage{amsmath,amssymb,amsthm}
\usepackage{mathtools}
\usepackage{hyperref}
\usepackage{doi}
\usepackage{geometry}
\usepackage{booktabs}
\usepackage{array}
\usepackage{enumitem}

\geometry{margin=2.5cm}

\input{glyphtounicode}
\pdfgentounicode=1

% The Adobe glyph list maps Delta and Omega to the technical codepoints U+2206 (INCREMENT) and
% U+2126 (OHM SIGN), so the deposited text layer would not be searchable as Greek letters.
\pdfglyphtounicode{Delta}{0394}
\pdfglyphtounicode{Omega}{03A9}
\pdfglyphtounicode{mu}{03BC}
\pdfglyphtounicode{periodcentered}{22C5}

\hypersetup{
  colorlinks=true,
  linkcolor=black,
  citecolor=black,
  urlcolor=black,
  pdftitle={The Gauge Structure Sub-Programme: Presentation Note 3},
  pdfauthor={J\'{e}r\^{o}me Beau},
  pdfsubject={Status of the gauge-sector identification after the O31 withdrawal: a structural
    U(1), a conditional SU(2), a withdrawn SU(3), and no established composite gauge group},
  pdfkeywords={gauge structure, admissible projection, U(1), SU(2), Hopf fibration,
    charge quantisation, withdrawn gauge-group derivation, Weil representation, presentation note}
}

\theoremstyle{definition}
\newtheorem{remark}{Remark}

\newcommand{\Heis}{\mathrm{Heis}}
\newcommand{\SU}{\mathrm{SU}}
\newcommand{\SO}{\mathrm{SO}}
\newcommand{\su}{\mathfrak{su}}
\newcommand{\UU}{\mathrm{U}}
\newcommand{\Gpi}{G_{\Pi}}
\newcommand{\Sym}{\mathrm{Sym}}
\newcommand{\ImH}{\operatorname{Im}\mathbb{H}}
\newcommand{\Ceff}{C_{\mathrm{eff}}}
\renewcommand{\Im}{\operatorname{Im}}

\title{The Gauge Structure Sub-Programme\\[0.4em]
\large Presentation Note 3}

\author{J\'{e}r\^{o}me Beau\\[0.3em]
\small Independent Researcher\\
\small \texttt{jerome.beau@cosmochrony.org}}

\date{2026\\[2pt]\small Version 2.3}

\begin{document}

  \maketitle

  \begin{abstract}
    The gauge structure sub-programme asks which internal symmetry groups are forced by the
    admissibility structure of the non-injective projection $\Pi$, rather than postulated.
    This note records what that question has and has not established.

    No derivation of the composite $\Gpi = \SU(3) \times \SU(2) \times \UU(1)$ is available.
    O31 withdraws every derivation of a gauge factor and of the Standard Model
    gauge group as a composite, together with the determinant-one condition, the commutation of the
    putative factors and the direct-product structure.
    The $\SU(3)$ identification is withdrawn with it: the uniqueness of $\SU(3)$ among connected
    compact subgroups acting irreducibly on $\mathbb{C}^3$ fails on the counterexample
    $\SO(3) \subset \SU(3)$; the irreducibility of the admissible transformation group fails on the
    diagonal torus $T^2$; and the arithmetic criterion $q \equiv 1 \pmod 3$ fails on
    $\{1, 2, q-3\}$.

    What remains is one structural sector and one conditional sector, not a group.
    $\UU(1)$ is structural: the Hopf phase fibre of $\Pi$ gives charge quantisation
    $q_{\mathrm{eff}} = w \cdot e$ from $\pi_1(S^1) \cong \mathbb{Z}$, on two premises carried as
    structural input rather than derived; and on the canonical admissible model no isolated magnetic
    monopole is topologically admissible.
    $\SU(2)$ is conditional on the spinor carrier supplied by O23: O27 gives
    unique real-linear factorisation through the admissibility quotient, and the identification of
    that quotient with $\su(2)$ is supplied rather than derived.

    \emph{Interpretation, not result}: the sub-programme's remaining content is a phase sector and a
    conditional non-abelian sector, and the step that would make them factors of a single gauge
    group is exactly the step that has been withdrawn.
  \end{abstract}

  \medskip
  \noindent\textbf{Keywords.}
  gauge structure; admissible projection; U(1); SU(2); Hopf fibration; charge quantisation;
  withdrawn gauge-group derivation; Weil representation; presentation note.

  \newpage
  \tableofcontents

  \newpage

  \section{Motivation and Central Question}
    \label{sec:motivation}

    The gauge structure sub-programme asks the following question.

    \medskip
    \noindent\textit{Central question.}
    The non-injective projection $\Pi: \chi \to O$ selects admissible configurations from the
    relational substrate.
    Its fibre $\Pi^{-1}(O_n)$ carries internal structure.
    \textit{Which internal symmetry groups are preserved as invariants of admissible projection?}

    \medskip
    The question is stated precisely as follows.
    Define $\Gpi$ as the group of fibre symmetries of $\Pi$ that preserve the admissibility
    constraints.
    Whether $\Gpi = \SU(3) \times \SU(2) \times \UU(1)$ is \emph{open}.
    The derivation of that identification is withdrawn at its source (O31~\cite{Beau2026a35}), and
    no replacement is available.

    \medskip
    Had the identification held, its physical importance would have been to give the Standard Model
    gauge group as a forced output of admissibility rather than a postulate.
    The framework's reading of gauge invariance as the residual fibre symmetry of a non-injective
    projection is unaffected, since it is a reading of what gauge freedom \emph{is} and not a
    derivation of which group appears.
    What is not available is the group.
    The spectral-entropy route to gauge dynamics (Q12~\cite{Beau2026q12},
    Q13~\cite{Beau2026q13}) does not take $\Gpi$ as a derived input: it carries $\Gpi$ as a
    supplied compact Lie group, and its results hold for every such group.

    \medskip
    This note is concerned exclusively with the \emph{identification} of the gauge group
    $\Gpi$.
    The derivation of Yang--Mills equations from the vertical variation of the spectral
    entropy belongs to the gauge--gravity spectral stratification sub-programme
    (Presentation Note~8) and is not claimed here.

  \section{Position in the Cosmochrony Programme}
    \label{sec:position}

    The gauge structure sub-programme sits at the intersection of Branch~I (axiomatic
    primitive and fibre structure) and Branch~II (spectral admissibility).
    It takes the following inputs from Presentation Notes~1 and~2 and uses them to address the
    identification of the internal symmetry group.

    From Note~1 (spectral admissibility):
    \begin{itemize}
      \item Spin-$\tfrac{1}{2}$ admissible sector $V_\rho \cong \mathbb{C}^2$, a supplied
      carrier (O23~\cite{Beau2026a27}, used by O27--O29~\cite{Beau2026a31,Beau2026a32,Beau2026a33};
      its selection from admissibility is open).
      \item Threshold dimension $\Sigma_c(n_3) = 3$, a supplied selection rule (O23
      Theorem~3.1~\cite{Beau2026a27} proves the three-dimensionality of the neutral sector
      of a supplied spinor carrier; the carrier selection and the identification of
      $\Sigma_c$ with that dimension are open).
      \item $\ImH \cong \su(2)$ as the three-dimensional adjoint sector of that carrier
      (O23--O27); its reading as the spatial algebra is a bridge of Note~2, not an input.
      \item No $\SU(3)$ colour-sector input. The single-frequency fingerprint mechanism that
      supplied it is withdrawn (O31~\cite{Beau2026a35}), and O32~\cite{Beau2026a36} supplies
      inter-triplet variance measurements, not the identification they were read as supporting.
    \end{itemize}

    From Note~2 (emergent geometry):
    \begin{itemize}
      \item Effective projection space $H_{\mathrm{eff}}$, identified with $\Sym^2(V_\rho)$ by a
      supplied hypothesis whose bridge requirements Q7~\cite{Beau2026q7} states.
      \item Effective base manifold $\mathbb{R}_\tau \times \Heis_3(\mathbb{R})$ (Q5b--Q11),
      with no metric coefficient supplied.
      Under the unestablished spatial-limit hypothesis [H-L] and the open lifting hypothesis
      [H-lift], Q5b~\cite{Beau2026q5b} gives a co-metric of rank three,
      $\mathrm{diag}(-A_\tau, A_H, A_H, 0)$, whose temporal term is supplied, with no coefficient
      values.
      Q11~\cite{Beau2026q11} gives a rank-four form $\mathrm{diag}(-A_\tau, 2\lambda, 2\lambda,
      2\lambda)$ only for a supplied full-rank extension under the hypotheses of its
      Corollary~6.1, with $A_\tau$ and $\lambda$ independent; it is proportional to $\eta$ only
      after a rescaling of $\tau$ (Note~2~\cite{Beau2026egnote}).
      The closure $g^{\mu\nu} = 2\eta^{\mu\nu}$ is not established, and nothing in this note uses
      a metric coefficient.
      \item Admissible principal bundle $P_{\Gpi}(M, \Gpi)$ for a supplied compact structure
      group $\Gpi$ (Q12~\cite{Beau2026q12}), with a supplied connection read as the generator of
      parallel transport, under two continuum hypotheses that do not determine it
      (Q6a~\cite{Beau2026q6a}).
    \end{itemize}

    \begin{remark}[Vocabulary note]
Throughout this note, \emph{gauge invariance} refers to the residual fibre symmetry of the
admissible projection $\Pi$, not to an independently postulated symmetry principle.
The substrate $\chi$ is static; gauge freedom arises from the non-injectivity of $\Pi$,
not from a dynamics of $\chi$.
    \end{remark}

    \begin{remark}[Scope: gauge group identification only]
This sub-programme concerns the identification of $\Gpi$, which it does not achieve.
It does not derive the Yang--Mills equations or the gauge--gravity coupling ratio either; those
belong to Note~8 (Gauge--Gravity Stratification Sub-Programme), whose results are stated for any
supplied compact $\Gpi$ and not for an identified Standard Model group, so the identification left
open here is not one of their inputs.
    \end{remark}

  \section{Logical Chain of the Sub-Programme}
    \label{sec:chain}

    The chain runs three stages and stops.
    Its fourth stage is withdrawn at the source, and the composite it would have terminated in is
    not established.

    \begin{equation}
      \label{eq:chain}
      \begin{aligned}
        &\underbrace{\Pi \simeq S^3}_{\text{fibre topology}}
        \;\Longrightarrow\;
        \underbrace{S^1 \hookrightarrow S^3 \to S^2}_{\text{Hopf fibration}}
        \;\Longrightarrow\;
        \underbrace{\UU(1)}_{\pi_1(S^1) \cong \mathbb{Z}} \\[1.2em]
        &\underbrace{\su(2) \text{ from } \ImH}_{\text{supplied carrier, O23}}
        \;\Longrightarrow\;
        \underbrace{\SU(2)}_{\text{conditional}}
        \;\not\Longrightarrow\;
        \underbrace{\SU(3)}_{\text{withdrawn}}
      \end{aligned}
    \end{equation}

    The arrows are of unequal type, and the last one is absent rather than weak.

    \begin{enumerate}
      \item \textit{$\UU(1)$ from the Hopf phase fibre.}
      The projection fibre is taken to be $\Pi^{-1}(O_n) \simeq S^3$, carrying a canonical Hopf
      fibration $S^1 \hookrightarrow S^3 \to S^2$; the circle fibre $S^1$ supports a $\UU(1)$
      phase symmetry with winding number $w \in \pi_1(S^1) \cong \mathbb{Z}$.
      Two premises here are structural inputs and not results.
      The identification of the fibre with $S^3$ is imported from the white paper and carried by
      TopologicalInvariants~\cite{Beau2026k} without proof; the note's own §\ref{sec:inputs}
      works on the complex fibre of a supplied Heisenberg carrier, which Foundation~M does not
      select, and the two identifications are reconciled only through the supplied carrier
      $V_\rho \cong \mathbb{C}^2$, whose unit sphere is $S^3$.
      And the identification of the $S^1$ phase with \emph{electric charge} is likewise taken as
      structural input there, not re-derived.
      What TopologicalInvariants Proposition~3.1 then gives, on those inputs, is
      $q_{\mathrm{eff}} = w \cdot e$.
      The $\UU(1)$ sector is argued in the topological layer and not in O31, so the O31 withdrawal
      notice does not reach it; that is a statement about its independence from O31, not a claim
      that its premises are derived.

      \item \textit{$\su(2)$ under a supplied carrier.}
      If the admissible neutral sector is carried by an irreducible two-dimensional unitary
      representation $V_\rho \cong \mathbb{C}^2$, a supplied input, then its traceless sector is
      $\su(2) \cong \ImH$ with $\dim_\mathbb{R} \ImH = 3$ (O23~\cite{Beau2026a27}, Theorem~3.1).
      O27~\cite{Beau2026a31} shows that a morphism whose kernel contains
      $\ker \Pi_{\mathrm{adm}}$ factors uniquely, as a real-linear map, through the admissibility
      quotient, and that the \emph{complex} intertwiner space of the representation action is
      one-dimensional.
      Identifying that quotient with $\su(2)$ is the carrier supplied by O23 under an open bridge
      there, not a consequence of the admissibility constraints, and the vector lift into
      $V_\rho$ remains open.
      The $\SU(2)$ sector is therefore conditional throughout.

      \item \textit{$\SU(3)$: withdrawn.}
      The colour-triplet route is withdrawn at its source.
      O31~\cite{Beau2026a35} withdraws the uniqueness of $\SU(3)$, the irreducibility of the
      admissible transformation group, the co-admissibility of the individual profiles
      $\sigma_{c_i}$, the arithmetic criterion, the metaplectic identity as used, the
      $\Delta(27)$ spectral proposition, a mistyped rank test, and every derivation of a gauge
      factor.
      Three counterexamples are given there in full: $\SO(3) \subset \SU(3)$ acts irreducibly on
      $\mathbb{C}^3$ with determinant one, so those two conditions never single out $\SU(3)$; the
      diagonal torus $T^2 \subset \SU(3)$ satisfies every imposed constraint while acting
      reducibly; and $\{1, 2, q-3\}$ is an admissible colour triplet for every prime $q \geq 7$,
      so triplets do not require $q \equiv 1 \pmod 3$.
      Nothing positive about $\SU(3)$ is asserted in this note.
    \end{enumerate}

  \section{Constituent Papers}
    \label{sec:papers}

    \subsection{The projective fibre and the \texorpdfstring{$\UU(1)$}{U(1)} phase sector}
      \label{sec:u1}

      The $\UU(1)$ factor arises from the topological structure of the projection fibre.

      \textbf{Foundation~M}~\cite{Beau2026m} does not identify the projection fibre
      $\Pi^{-1}(O_n)$ with a Heisenberg or Weil carrier: its Theorem~5.7 shows that the carrier
      contract extracted from its axioms does not select $\Heis_3(\mathbb{Z}/q\mathbb{Z})$.
      The spectral programme supplies that carrier, and the phase degree of freedom used below is
      the phase of its complex fibre.

      \textbf{TopologicalInvariants}~\cite{Beau2026k} classifies the admissible homotopy
      invariants of the projection fibre.
      The projection fibre $\Pi \simeq S^3$ carries the Hopf fibration
      $S^1 \hookrightarrow S^3 \xrightarrow{\pi_H} S^2$.
      Three results are established.
      (i)~Electric charge arises from $\pi_1(S^1) \cong \mathbb{Z}$: the winding number
      $w \in \mathbb{Z}$ is the topological invariant, giving $q_{\mathrm{eff}} = w \cdot e$
      (Proposition~3.1).
      That proposition does not invoke the canonical admissible model; its undischarged inputs are
      the identification of the fibre with $S^3$ and the identification of the $S^1$ phase with
      electric charge, both carried as structural input there.
      (ii)~On the canonical admissible model, the bounded-flux constraint removes the polar
      regions of the Hopf base $S^2$, reducing the admissible base to an annulus with
      $\pi_2(\mathcal{C}_{\mathrm{adm}}) = 0$ (Proposition~4.1): no isolated magnetic monopoles
      are topologically admissible.
      (iii)~Parity violation is interpreted as projective chirality: an orientation dependence of
      admissibility under $\Pi$, which selects one handedness of winding invariants.

      \textbf{Q6a}~\cite{Beau2026q6a} defines the gauge group as the group of fibre-preserving
      transformations of $\Pi_n$.
      It does not determine the full gauge group, and it does not use O31: it treats $\SU(3)$ as
      a structural hypothesis whose identification is open.
      Its $\UU(1)$ is the group of global phase rotations $\omega \mapsto e^{i\theta}\omega$ of the
      complex fibre; they act as a global scalar, which leaves every span-based capacity profile
      unchanged, so the global phase is not resolved by those observables, and the rotations belong to
      the gauge group to the extent that $\Pi_n$ is defined through such observables.
      Q6a states two continuum hypotheses that it does not prove and shows that they do not determine
      a connection; on a supplied structure group, representation and connection, the gauge connection
      is the infinitesimal generator of parallel transport.

      \textbf{SMchargemass}~\cite{Beau2026e} uses the $\UU(1)$ gauge-structure result as a
      downstream input, not as a constituent of its identification.
      Mass, charge, and inertia are interpreted there as saturated-response regimes of the
      same bounded projective mechanism.
      It is therefore an application of the gauge-structure output established here, and is
      listed under Outputs to Downstream Branches rather than as a constituent paper.

    \subsection{Quaternionic admissibility and the \texorpdfstring{$\SU(2)$}{SU(2)} sector}
      \label{sec:su2}

      The $\SU(2)$ factor is studied by the chain O23--O30 and Q6a, conditionally on the
      supplied spinor carrier.

      \textbf{O23}~\cite{Beau2026a27} proves that if the admissible neutral sector is carried
      by an irreducible two-dimensional unitary representation $V_\rho \cong \mathbb{C}^2$ (a
      supplied input), then its traceless sector is $\su(2) \cong \ImH$ with
      $\dim_\mathbb{R} \ImH = 3$ (Theorem~3.1).
      The threshold $\Sigma_c(n_3) = 3$ is a supplied selection rule: the carrier selection
      and the identification of $\Sigma_c$ with that dimension are open bridges.

      \textbf{O27}~\cite{Beau2026a31} proves that a morphism whose kernel contains
      $\ker \Pi_{\mathrm{adm}}$ factors uniquely, as a real-linear map, through the admissibility
      quotient, and that under the supplied carrier the \emph{complex} intertwiner space of the
      representation action $\rho \colon \su(2) \to \mathrm{End}(V_\rho)$ is one-dimensional, the
      corresponding real-linear space having real dimension two (Theorem~4.1 there).
      Identifying the quotient with $\su(2)$ is the carrier supplied by O23, not a consequence of
      the admissibility constraints, and there is no non-zero $\su(2)$-equivariant map from the
      adjoint to the fundamental, so the vector lift into $V_\rho$ is open.
      This is a factorisation statement, not a rigidity theorem, and it is conditional.

      \textbf{O28}~\cite{Beau2026a32} reports, in its protocol, $1\%$ threshold rank 3 for the
      covariance $\mathcal{C}_c$ in $\mathrm{End}(H_{\mathrm{eff}})$ at $q \in \{61, 101, 151\}$, with
      eigenvalue structure $[1 : \tfrac{1}{2} : \tfrac{1}{2}]$.
      These are measured values of that protocol, not invariants.

      \textbf{O29}~\cite{Beau2026a33} audits that finite-data result: rank three is modal but not
      invariant ($32/50$ stored samples at $q=101$, $103/105$ at $q=211$).
      It separates the measured span in $H_{\mathrm{eff}}$, its data-dependent truncation, and the
      supplied spin-$\tfrac{1}{2}$ module $V_\rho$, and identifies none of them with another.
      The doublet dimension $d_\rho = 2$ is a minimality selection of O26 throughout this note, not
      a measurement.

      \textbf{O30}~\cite{Beau2026a34} studies a conditional $3\times3$ covariance model on
      $\Sym^2(V_\rho)$ under explicit Veronese-cone, inter-block and parity assumptions; it is not the
      $9\times9$ covariance measured in O28 and supplies no analytical account of the O28 ratio.
      It is not load-bearing for anything asserted below.

      \textbf{Q6a}~\cite{Beau2026q6a} treats the $\SU(2)$ sector at the level of the gauge group of the
      admissible fibre: on the spinor carrier supplied by O23, and under the neutral-sector
      hypothesis of O27, the group acting faithfully on the carrier is $\SU(2)$ (its
      Proposition~2.5), a consistency statement rather than a selection; its connection is supplied,
      and it generates parallel transport on supplied data under the same two continuum hypotheses as
      in §\ref{sec:charges}, which do not determine it.
      That construction presupposes the sector rather than closing it, so it inherits the
      conditionality of the supplied carrier.

      \textbf{Q7}~\cite{Beau2026q7} states what an equivariant bridge between the spatial and internal
      structures would require: under the hypothesis identifying the measured space with a spin-1
      module, $\Sym^2(V_\rho)$ is the unique irreducible three-dimensional $\su(2)$-module, so a
      non-zero equivariant bridge is an isomorphism.
      That identification is supplied, not derived.

    \subsection{Gauge charges and admissible bundle structure}
      \label{sec:charges}

      \textbf{Q6a}~\cite{Beau2026q6a} defines a gauge charge as a functional constant on the
      projection fibres of $\Pi_n$.
      Such a functional factors through $\Pi_n$, so it is a function of the observable state, and its
      identification with a conserved physical charge is open in Q6a.
      Its bundle-level content is conditional on two hypotheses that it does not prove: that the
      admissible fibre forms a smooth fibre bundle over an effective manifold, and that this
      manifold carries a differentiable structure.
      These hypotheses do not determine a connection, so Q6a supplies a structure group, its
      representation and a connection; on these data the connection generates parallel transport, and
      the field strength is its curvature, which measures the path dependence of transport.
      Q6a supplies no decomposition of the fibre into orbit sectors.
      The white paper shows that injective admissible projections carry trivial holonomy; Q6a's
      supplied transport is the local infinitesimal counterpart of that setting and establishes no
      non-trivial holonomy for a given projection.

      \textbf{Q12}~\cite{Beau2026q12} defines the admissible principal $\Gpi$-bundle for a supplied
      compact Lie group $\Gpi$ and extends $A_g$ to the Laplace-type operator $A_{g,A}$ on the
      associated vector bundle.
      It takes no specific group: its main theorem holds for any compact $\Gpi$, and no result of
      it depends on which group is supplied, on a decomposition of the bundle into sub-bundles, or
      on a colour sector.
      Its Lemma~1, titled there \emph{Vertical variation at fixed base metric}, states under
      minimal coupling --- the endomorphism $E$ not depending on the connection $A$ --- that the
      vertical variation at fixed $g$ receives no $a_2$ contribution.
      Q12 itself follows it with a remark titled \emph{No full horizontal--vertical decoupling}, so
      it is not the stronger independence statement its usual name suggests.

      \textbf{TopologicalInvariants}~\cite{Beau2026k} places the charge identification within
      the topological layer: charge is a $\pi_1$ invariant; the absence of monopoles follows
      from $\pi_2(\mathcal{C}_{\mathrm{adm}}) = 0$ on the canonical model.

    \subsection{The withdrawn \texorpdfstring{$\SU(3)$}{SU(3)} route}
      \label{sec:su3}

      The colour-triplet route to $\SU(3)$ is withdrawn at its source and supplies nothing to this
      sub-programme.

      \textbf{O31}~\cite{Beau2026a35} is a withdrawal notice.
      It withdraws eight results of its previous version: the uniqueness of $\SU(3)$ among
      connected compact subgroups acting irreducibly on $\mathbb{C}^3$; the irreducibility of the
      group of Born--Infeld-compatible fibre-preserving transformations; the co-admissibility of
      the individual profiles $\sigma_{c_1}, \sigma_{c_2}, \sigma_{c_3}$; the criterion that colour
      triplets exist if and only if $q \equiv 1 \pmod 3$; the metaplectic identity as used, which
      changes the central character while the automorphism invoked fixes the centre; the
      $\Delta(27)$ spectral proposition; a rank-eight test mistyped in $\mathrm{End}(\mathbb{C}^3)$,
      where the dimension is nine; and every derivation of a gauge factor and of
      $\SU(3) \times \SU(2) \times \UU(1)$ as a composite, the determinant-one condition, the
      commutation of the putative factors and the direct-product structure included.
      Its remaining scalar, arithmetic and discrete-group material, the generator-twist
      isospectrality statement among it, is under re-audit and is not asserted there.

      Three counterexamples make the withdrawal unavoidable rather than cautionary.
      $\SO(3) \subset \SU(3)$, the real orthogonal matrices of determinant one read as complex
      matrices, is a proper connected compact subgroup whose action on $\mathbb{C}^3$ is
      irreducible, so irreducibility together with determinant one never singles out $\SU(3)$.
      The diagonal torus $T^2 \subset \SU(3)$ satisfies every constraint that was imposed while
      preserving the three coordinate lines, so it acts reducibly; requiring each sector to be
      mapped to itself in fact forces reducibility.
      And $\{1, 2, q-3\}$ is an admissible colour triplet under the published definition for every
      prime $q \geq 7$, so the arithmetic criterion is false; $q = 11$ admits $\{1, 2, 8\}$ with
      $11 \equiv 2 \pmod 3$.

      \textbf{O32}~\cite{Beau2026a36} remains a numerical campaign.
      It is cited here for what it measures: inter-triplet variance ratios $R_{\mathrm{var}}$ at or
      below a control-calibrated noise floor for $q \in \{61, 151, 211, 307\}$, against controls
      $q \in \{29, 101\}$.
      Its Proposition~1, a modulation-frequency argument giving
      $\mathbb{E}[\sigma_c(n)] = \mathbb{E}[\sigma_{\omega c}(n)] + O(q^{-1})$, is an O32 result
      that O31 does not withdraw; it is cited as published and is not affected by anything here.
      Those results stand.
      What does not stand is the reading of them as confirming a hypothesis whose consequence,
      the $\SU(3)$ identification, has been withdrawn: a variance ratio at the noise floor is
      evidence about profile equality, not about a symmetry group.

    \subsection{Boundary with Yang--Mills dynamics}
      \label{sec:ym-boundary}

      The gauge structure sub-programme concerns the identification of $\Gpi$, which it does not
      achieve, and records the admissible principal bundle of a supplied compact structure group.
      It does not derive the Yang--Mills equations either.
      That derivation belongs to the gauge--gravity spectral stratification sub-programme
      (Presentation Note~8), where the vertical variation of the joint spectral entropy
      functional $\mathcal{S}_\Pi[g, A]$ yields the $a_4$ Yang--Mills response.

      The logical boundary is sharp: Q6a and Q12 are constituent papers of the gauge
      structure sub-programme in so far as they address the gauge group and construct
      the bundle.
      Q12 and Q13 become constituent papers of Note~8 in so far as they derive the
      Yang--Mills equations and the gauge--gravity coupling.

      The $a_4$-level result of Q12 (Theorem~1) takes $\Gpi$ as a supplied compact group and holds
      for every such group; the group enters only through the fibre representation and its Dynkin
      index.
      It therefore does not depend on the identification that this sub-programme leaves open, and
      no result of Q12 depends on a colour sector.

      \begin{table}[ht]
        \centering
        \caption{Summary of constituent papers by sector.
        Status: P = proved; S = structural; N = numerical; C = conditional; O = open;
        W = withdrawn at the source.}
        \label{tab:papers}
        \renewcommand{\arraystretch}{1.2}
        \begin{tabular}{@{}llll@{}}
          \toprule
          Sector & Papers & Central output & Status \\
          \midrule
          $\UU(1)$ fibre & TopInv & Charge $= w \cdot e$, on imported inputs & S \\
          Monopole exclusion & TopInv & $\pi_2(\mathcal{C}_{\mathrm{adm}}) = 0$, canonical model & P \\
          $\UU(1)$ phase & Q6a & Global phase rotations, unresolved by span-based observables & S \\
          $\SU(2)$ sector & O23, O27, Q6a, Q7 & Quotient factorisation; $\su(2)$ supplied & C \\
          Charges & Q6a, TopInv & Fibre-invariant charges & S \\
          Bundle & Q12, Q6a & Supplied compact $\Gpi$ and connection & S, C \\
          $\SU(3)$, either graph & O31 & Triplet co-admissibility route & W \\
          Triplet variance ratios & O32 & $R_{\mathrm{var}}$ at the noise floor & N \\
          $\Gpi$ as a composite & O31 & Direct-product structure & W \\
          \bottomrule
        \end{tabular}
      \end{table}

  \section{Inputs from Upstream Results}
    \label{sec:inputs}

    \begin{enumerate}
      \item \textbf{Supplied carrier and fibre structure.}
      The identification of $\Pi^{-1}(O_n)$ with a Heisenberg carrier is supplied, not a theorem:
      Foundation~M~\cite{Beau2026m} shows that its axioms do not select it (its Theorem~5.7).
      On that supplied carrier, the phase degree of freedom is the phase of the complex fibre.

      \item \textbf{Non-injectivity as structural necessity.}
      ENI~\cite{Beau2026l} establishes that non-injectivity is not an interpretive option:
      any genuinely emergent description with $S_\Pi > 0$ must be non-injective.
      Gauge freedom is a direct consequence of this non-injectivity: the fibre
      $\Pi^{-1}(O_n)$ is non-trivial precisely because $\Pi$ is non-injective.

      \item \textbf{Spin-$\tfrac{1}{2}$ sector and $\Sigma_c(n_3) = 3$.}
      From Note~1 (O23, O27, O29), as supplied inputs: the admissible sector is taken to be
      carried by $V_\rho \cong \mathbb{C}^2$ (supplied carrier, not derived);
      $\ImH \cong \su(2)$ is the three-dimensional adjoint sector of that carrier (O23,
      Theorem~3.1); the threshold $\Sigma_c(n_3) = 3$ and the minimality selection $d_\rho = 2$ are
      supplied selection rules, their derivation from admissibility being open.

      \item \textbf{Effective geometry.}
      From Note~2 (Q5b--Q11): the effective base manifold $\mathbb{R}_\tau \times
      \Heis_3(\mathbb{R})$.
      Q12 defines the principal bundle over a Riemannian four-manifold $M$ with a supplied metric,
      and reading $M$ as this emergent base is conditional.
    \end{enumerate}

  \section{Outputs to Downstream Branches}
    \label{sec:outputs}

    \begin{table}[ht]
      \centering
      \caption{Principal outputs and downstream consumers.
      ---: no consuming paper is identified.
      Q12 and Q13 take $\Gpi$ as a supplied compact group, and Q14~\cite{Beau2026q14} carries
      its colour module as an explicit hypothesis; none of them consumes a withdrawn output.}
      \label{tab:outputs}
      \renewcommand{\arraystretch}{1.3}
      \begin{tabular}{@{}p{5.5cm}p{4.5cm}l@{}}
        \toprule
        Output & Consumer & Status \\
        \midrule
        $\Gpi = \SU(3) \times \SU(2) \times \UU(1)$ & --- & W (withdrawn, O31) \\
        Charge/mass/inertia as saturated responses & SMchargemass~\cite{Beau2026e} & S \\
        Admissible principal bundle $P_{\Gpi}(M, \Gpi)$ & Q13, Q14 & S, for a supplied compact $\Gpi$ \\
        Charge quantisation $q_{\mathrm{eff}} = w \cdot e$ & --- & S (imported inputs) \\
        Absence of magnetic monopoles & --- & P (canonical) \\
        $\su(2)$ sector under a supplied carrier & --- & C \\
        $\SU(3)$ colour sector & --- & W (withdrawn, O31) \\
        Projective chirality & --- & S \\
        \bottomrule
      \end{tabular}
    \end{table}

  \section{Status of Results}
    \label{sec:status}

    \subsection{Proved results}
      \label{sec:proved}

      \begin{enumerate}
        \item \textit{Absence of magnetic monopoles on the canonical model}
        (TopologicalInvariants, Proposition~4.1 and Corollary~5.1):
        $\pi_2(\mathcal{C}_{\mathrm{adm}}) = 0$ on the canonical admissible model;
        no isolated magnetic charge is topologically admissible.

        \item \textit{Neutral-sector dimension under a supplied carrier}
        (O23~\cite{Beau2026a27}, Theorem~3.1):
        the traceless sector of the supplied spinor carrier is $\su(2) \cong \ImH$ with
        $\dim_\mathbb{R} \ImH = 3$.
        The carrier selection, and the identification of $\Sigma_c(n_3) = 3$ with that dimension,
        are open bridges of O23.

        \item \textit{Vertical variation at fixed base metric} (Q12~\cite{Beau2026q12}, Lemma~1):
        under minimal coupling, the vertical variation at fixed $g$ receives no $a_2$ contribution.
        Q12 states separately that full horizontal--vertical decoupling does \emph{not} hold.
        The lemma does not depend on which group occupies the gauge sector, which is why it survives
        the withdrawal.
      \end{enumerate}

    \subsection{Conditional results}
      \label{sec:conditional}

      \begin{enumerate}
        \item \textit{Factorisation through the admissibility quotient}
        (O27~\cite{Beau2026a31}, Theorem~4.1):
        a morphism whose kernel contains $\ker \Pi_{\mathrm{adm}}$ factors uniquely, as a
        real-linear map, through the admissibility quotient, and the \emph{complex} intertwiner
        space of the representation action is one-dimensional under the supplied carrier.
        Identifying the quotient with $\su(2)$ is supplied by O23 and rests, in O27, on the
        admissible-saturation hypothesis, not derived; O27 carries a second hypothesis, parity
        transport, for the involution's descent.  The vector lift into $V_\rho$ is open.

        \item \textit{$\SU(2)$ gauge sector} (Q6a~\cite{Beau2026q6a}):
        the $\SU(2)$ factor is read off the supplied spinor carrier, under O27's neutral-sector
        hypothesis (Q6a Proposition~2.5), with the
        gauge field a supplied connection, the generator of parallel transport under Q6a's two
        continuum hypotheses, which do not determine it.
        The construction presupposes the sector and inherits the carrier's conditionality.
      \end{enumerate}

    \subsection{Structural results}
      \label{sec:structural}

      \begin{enumerate}
        \item \textit{Charge quantisation}
        (TopologicalInvariants~\cite{Beau2026k}, Proposition~3.1):
        $q_{\mathrm{eff}} = w \cdot e$ with $w \in \mathbb{Z} = \pi_1(S^1)$, on two structural
        inputs that TopologicalInvariants carries rather than derives: the identification of the
        projection fibre with $S^3$, and the identification of the $S^1$ phase with electric charge.
        Proposition~3.1 does not invoke the canonical admissible model; that restriction belongs to
        the monopole results below.

        \item \textit{$\UU(1)$ as phase symmetry} (Foundation, Q6a):
        the $\UU(1)$ sector is the group of global phase rotations of the complex fibre, which
        leave every span-based capacity profile unchanged (Q6a); the fibre is that of a supplied
        Heisenberg carrier, which Foundation~M does not select.

        \item \textit{Gauge charges as fibre invariants} (Q6a): a gauge charge is a functional
        constant on the projection fibres, hence a function of the observable state, whose
        identification with a conserved physical charge is open.
        The connection and its curvature $F_{\mu\nu}$ are defined only on supplied transport data,
        under Q6a's two continuum hypotheses, which do not determine them; Q6a does not determine in
        full the gauge group of fibre-preserving transformations.

        \item \textit{Projective chirality and parity} (TopologicalInvariants, §6):
        parity violation arises from the orientation dependence of admissibility under $\Pi$;
        the $V - A$ structure is identified as the maximal chirality class
        (partial structural result).
      \end{enumerate}

    \subsection{Numerical evidence}
      \label{sec:numerical}

      \begin{enumerate}
        \item \textit{$1\%$ threshold rank 3} (O28--O29):
        measured at $q \in \{61, 101, 151\}$ in O28's protocol; in O29's extended recomputation rank
        three is modal, not invariant.
        The doublet dimension $d_\rho = 2$ is a minimality selection of O26 (§\ref{sec:inputs}),
        not a measurement.

        \item \textit{Inter-triplet variance ratios} (O32~\cite{Beau2026a36}):
        $R_{\mathrm{var}}$ at or below a control-calibrated noise floor for
        $q \in \{61, 151, 211, 307\}$, controls $q \in \{29, 101\}$.
        Evidence about profile equality, not about a symmetry group.
      \end{enumerate}

    \subsection{Withdrawn}
      \label{sec:withdrawn}

      The following are withdrawn on the authority of the withdrawal notice
      O31~\cite{Beau2026a35}, and are therefore not asserted here.

      \begin{enumerate}
        \item \textit{The $\SU(3)$ identification}, on the colour-adapted graph and on the
        standard graph alike, at every analytical level.
        \item \textit{Triplet co-admissibility} as a proved consequence of metaplectic
        intertwining.
        \item \textit{The arithmetic criterion} that colour triplets require $q \equiv 1 \pmod 3$.
        \item \textit{Generator-twist isospectrality} as a proved closure of the spectral route;
        it is under re-audit at the source and is not asserted.
        \item \textit{$[\mathrm{H\text{-}color}]$ at the rank level} via intra-coset Vandermonde
        symmetry, which belongs to the material under re-audit.
        \item \textit{The composite} $\Gpi = \SU(3) \times \SU(2) \times \UU(1)$, together with
        the determinant-one condition, the commutation of the putative factors and the
        direct-product structure.
      \end{enumerate}

    \subsection{Open hypotheses}
      \label{sec:open-hyp}

      \begin{enumerate}
        \item \textit{Any route to $\SU(3)$.} No mechanism is currently available.
        The colour-triplet route is withdrawn, and the counterexamples of O31 close the specific
        arguments rather than merely leaving them unfinished.

        \item \textit{Any route to a composite $\Gpi$.} Even granting $\UU(1)$ and a conditional
        $\SU(2)$, nothing establishes that they are factors of a single group, that the factors
        commute, or that the structure is a direct product.

        \item \textit{$\pi_2(\Ceff) = 0$ on the full configuration space}
        (TopologicalInvariants, Remark~4.2):
        the monopole exclusion is proved on the canonical model; extension to the full effective
        configuration space $\Ceff$ is open.

        \item \textit{Maximal chirality ($V - A$)} (TopologicalInvariants, Remark~6.2):
        projective chirality is established as a structural mechanism; its identification with the
        specific $V - A$ structure of the weak interaction remains qualitative.
      \end{enumerate}

  \section{Open Deliverables}
    \label{sec:open}

    \textbf{Open deliverable 1: a route to $\SU(3)$.}
    The colour-triplet route is withdrawn and its specific arguments are refuted by
    counterexample, not merely left incomplete.
    A replacement would have to supply a mechanism that singles out $\SU(3)$ among the connected
    compact subgroups acting irreducibly on $\mathbb{C}^3$, which irreducibility and determinant
    one demonstrably do not, and it would have to do so without the withdrawn metaplectic step.

    \textbf{Open deliverable 2: a route to a composite gauge group.}
    Suppose $\UU(1)$ and a conditional $\SU(2)$ are granted.
    Nothing establishes that they are factors of one group acting on one fibre, that they commute,
    or that the composite is a direct product.
    O31 withdrew each of those three separately, and this note asserts none of them.

    \textbf{Open deliverable 3: $\pi_2(\Ceff) = 0$ in general.}
    The topological exclusion of magnetic monopoles is proved on the canonical admissible model.
    Extension to the full $\Ceff$ is a well-posed open problem in homotopy theory.

    \textbf{Open deliverable 4: $V - A$ maximality from admissibility.}
    The structural mechanism for projective chirality is identified.
    Deriving the specific $V - A$ form of the weak interaction current from admissibility data is
    an open analytical problem.

  \section{Conclusion}
    \label{sec:conclusion}

    The gauge structure sub-programme does not identify an admissible internal symmetry group.

    Of its three sectors, one is structural, one is conditional and one is withdrawn.
    $\UU(1)$ arises from the Hopf phase fibre of $\Pi \simeq S^3$, with charge quantisation
    following from $\pi_1(S^1) \cong \mathbb{Z}$ on two premises carried as structural input; the
    monopole exclusion is the result that holds on the canonical admissible model.
    $\SU(2)$ arises from the spinor carrier supplied by O23, whose neutral sector is
    $\su(2) \cong \ImH$; O27 supplies factorisation through the admissibility quotient and a
    one-dimensional \emph{complex} intertwiner space, while the identification of that quotient with
    $\su(2)$
    remains supplied rather than derived.
    $\SU(3)$ has no surviving route: O31 withdraws the uniqueness, the irreducibility, the
    co-admissibility of the individual profiles, the arithmetic criterion, the metaplectic step and
    every derivation of a gauge factor, with three counterexamples given in full.

    The composite $\Gpi = \SU(3) \times \SU(2) \times \UU(1)$ is therefore not established, and
    neither are the determinant-one condition, the commutation of the putative factors, nor the
    direct-product structure.
    Q12 and Q13 carry $\Gpi$ as a supplied compact group and hold for every such group, and Q14
    carries its colour module as an explicit hypothesis; none of them takes $\Gpi$ as a derived
    input, so their results are statements about a supplied group, not about an identified
    Standard Model group.

    This note does not derive Yang--Mills dynamics, which belongs to Note~8
    (Gauge--Gravity Stratification).
    Note~8 takes $\Gpi$ as a supplied compact group.

    \subsection*{Interpretive outlook}

    What follows is a structural reading, offered as interpretation and not as a result.

    The ambition of a gauge sub-programme is that gauge structure be not postulated but forced:
    that the Standard Model group fall out of admissibility.
    What is available is weaker and more specific.
    Non-injectivity does give a non-trivial fibre, and a fibre does carry a phase symmetry, so the
    reading of gauge freedom as residual fibre symmetry is intact as a reading.
    What is not available is the inference from that reading to any particular group.
    The withdrawn arguments failed in a characteristic way: they imposed conditions that a target
    group satisfies and concluded that the target was thereby singled out, which
    $\SO(3) \subset \SU(3)$ and the diagonal torus each refute directly.
    A future route will have to identify a group by an argument that excludes its subgroups, not
    merely by one that the group passes.

    \newpage
    \bibliographystyle{plain}
    \bibliography{cosmochrony-bibliography}

\end{document}
